\section{Conclusion}
\label{sec:conclusion}

\paragraph{The answer to the guiding question.}
The question of Section~\ref{subsec:question} was why the inviscid Burgers'
equation needs an entropy condition and why that condition is the right one.
The chain of results is now complete. Smooth decreasing data develop an
infinite slope at an explicit time and place,
$T^*=-1/\min u_0'$ (Section~\ref{sec:shock-formation}); the weak formulation
continues the solution past this point through jumps obeying the
Rankine--Hugoniot condition (Section~\ref{sec:weak-solutions}), but it is
non-unique, already for the simplest Riemann data, where in fact infinitely
many weak solutions exist (Section~\ref{sec:nonuniqueness}). The Oleinik
entropy condition, of which the Lax condition $u_\ell>u_r$ is the local form
at a jump, restores uniqueness and gives the complete solution of the Riemann
problem (Sections~\ref{sec:entropy-condition}--\ref{sec:riemann-problem}).
That it is the \emph{right} condition is the content of the vanishing
viscosity limit: via the Cole--Hopf transformation the viscous solutions
converge to the entropy solution, whose jumps can only be downward because
the Lax--Oleinik minimiser is monotone, and directly because a viscous
travelling wave connecting $u_\ell$ to $u_r$ exists exactly when
$u_\ell>u_r$ (Section~\ref{sec:extensions}, Proposition~\ref{prop:viscous-profile}). The
theory itself is standard \cite[\S3.4]{evans2010}; the novelty of this
paper is in the explicit verifications and the computations below.

\paragraph{What the computations added.}
\begin{itemize}
  \item \emph{An explicit example.} For $u_0=-\arctan x$ the entropy solution
    after breakdown is given in closed parametric form, with a single
    stationary shock at $x=0$ of strength $\arctan y^*(t)$, where $y^*$ is
    the positive root of $\arctan y=y/t$. The strength grows like
    $\sqrt{3(t-1)}$ from the breakdown time $t=1$ and saturates at $\pi/2$
    (Section~\ref{subsec:arctan-complete}).
  \item \emph{A prediction that was confirmed.} A viscous shock has $L^1$
    distance exactly $4\ln2\,\nu\approx2.77\,\nu$ from the sharp shock,
    independent of its strength and speed. In the Cole--Hopf solutions the
    shock layer contributes $2.78\,\nu$, within $0.5\%$ of this, and the
    contribution tends to $4\ln2\,\nu$ at every strength tested
    ($0.52$--$1.50$); it is visibly larger in a wide window only for young
    shocks, where the steep outer solution contaminates it. The total
    $L^1$ error is first order in $\nu$, tending to $\approx3.11\,\nu$,
    the remaining $\approx0.33\,\nu$ coming from the smooth part (which is
    measured, not derived).
  \item \emph{Godunov as a vanishing viscosity.} The scheme converges at
    first order for the shock, with error exactly $0.50\,\Delta x$ when the
    shock lies on a cell interface and $0.6$--$0.77\,\Delta x$ for the misaligned grids tested;
    for the rarefaction the observed order is about $0.8$ and still rising,
    so a first-order rate is not established there. The modified equation
    predicts the shock error to within $3\%$ at CFL $0.45$ (and within about
    $10\%$ for CFL numbers from $0.1$ to $0.95$), corresponding to an
    effective viscosity $\approx0.18\,\Delta x\to0$: a monotone scheme
    selects the entropy solution because its numerical diffusion acts as a
    vanishing viscosity. This is a heuristic consistency check; the rigorous
    result is the convergence theorem for monotone schemes
    \cite{crandall1980}.
  \item \emph{Corrected N-wave constants.} The long-time N-wave is
    determined by $Q=-\min V$ and $P=\int u_0-\min V$, where $V$ is the
    primitive of $u_0$, \emph{not} by the positive and negative masses of
    $u_0$: positive and negative parts cancel across shocks, and only
    $P-Q=\int u_0$ is conserved. For the data $+1$ on $(0,1)$, $-1$ on
    $(1,2)$ the correct limit is $0$ (the $L^1$ norm decays exactly like
    $1/t$), while the positive/negative masses would predict a wave of
    norm $2$. Numerically the distance to the N-wave with the correct
    constants decays in every case tested, and with the naive constants it
    does not in the cases where the two differ (Section~\ref{subsec:n-wave}).
\end{itemize}

\paragraph{Limits of the exposition.}
The proofs of the uniqueness theorem (Theorem~\ref{thm:oleinik-uniqueness}),
of the vanishing-viscosity limit (Theorem~\ref{thm:vanishing-viscosity}), of
the convergence of the Godunov scheme (Theorem~\ref{thm:godunov-convergence})
and of the N-wave decay (Theorem~\ref{thm:n-wave-decay}) are cited, not
reproduced. The numerical statements are observations for the specific data
sets and resolutions stated. Burgers' equation is a model of nonlinear
steepening and its regularisation, not a stepping stone to the far harder
open regularity questions for the three-dimensional Euler and Navier--Stokes
equations, whose mechanisms differ.